% Fichier généré automatiquement par tools/generate_corrections.py.
% Source : DYN/DYN-04-TorseurDynamique/1025_RTR/1025_RTR.tex
% Statut : complété (4/4 question(s)).
\begin{corrigebox}[Corrigé complété]
\CorrigeQuestion{1}
\textbf{Calculons $\vectrd{2}{0}$}

$\vectrd{2}{0} = m_2\vectg{G_2}{2}{0}=m_2 \vectg{D}{2}{0} $

 \textbf{Calcul de $\vectv{D}{2}{0}$ : }  
$\vectv{D}{2}{0}=\deriv{\vect{AD}}{\rep{0}}$ 
$=\deriv{ R\vi{1} + \lambda (t) \vi{1}}{\rep{0}}$  
$= \left( R + \lambda (t)\right) \thetap \vj{1} + \lambdap (t) \vi{1}$.

 \textbf{Calcul de $\vectg{D}{2}{0}$ : }  
$\vectg{D}{2}{0}=\deriv{\vectv{D}{2}{0}}{\rep{0}}$ 
$=\deriv{  \left( R + \lambda (t)\right) \thetap \vj{1} + \lambdap (t) \vi{1}}{\rep{0}}$  
$= \left(\lambdap (t)\right) \thetap \vj{1} +\left( R + \lambda (t)\right) \thetapp \vj{1} - \left( R + \lambda (t)\right) \thetap^2 \vi{1} + \lambdap (t) \vi{1}$
\CorrigeQuestion{2}
On choisit les inconnues et les conventions positives de la figure,
        on écrit les relations de conservation/fermeture adaptées, puis on résout le
        système obtenu. Le résultat symbolique doit être homogène et vérifié dans les cas
        limites (entrée nulle, rapport unitaire ou position de référence).
\CorrigeQuestion{3}
On choisit les inconnues et les conventions positives de la figure,
        on écrit les relations de conservation/fermeture adaptées, puis on résout le
        système obtenu. Le résultat symbolique doit être homogène et vérifié dans les cas
        limites (entrée nulle, rapport unitaire ou position de référence).
\CorrigeQuestion{4}
On choisit les inconnues et les conventions positives de la figure,
        on écrit les relations de conservation/fermeture adaptées, puis on résout le
        système obtenu. Le résultat symbolique doit être homogène et vérifié dans les cas
        limites (entrée nulle, rapport unitaire ou position de référence).
\end{corrigebox}
